\section{Evaluation Protocol}
\label{sec:evaluation}

\subsection{Metrics}

The main outcome is end-to-end task success: the final artifact satisfies
the task predicate within the budget and without violating the sandbox policy.
Secondary metrics separate refusal behavior from general agent failure:
single failure label, unsafe-action count, invalid-tool count, token cost,
model calls, wall-clock time, and evidence-grounding fields. Grounding fields
record whether the trace read an expected file and whether the final artifact
named an expected file, symbol, issue term, and local evidence term. This
distinction matters because a lower success rate is not evidence of
alignment-induced suppression unless the trace shows refusal, suppressed
action, or phase-specific omission rather than generic tool or capability
failure. The metrics are defined in
Appendix~\ref{app:supplementary-tables}.

\subsection{Sample Size}

The 31B Gemma repair-controller security evaluation has 30 tasks, five seeds,
and four conditions, for 600 runs. The 26B Gemma, Qwen2.5-Coder, and Llama 3.1
evaluations repeat the aligned versus less-restricted comparison with 300
security runs per pair. Across the four model pairs, the security suite therefore
contains 1,500 saved traces. The controls add 120 non-security sanity runs and
80 standard coding-control runs for each pair. Separate 24-run pilots validate
orchestration; a 150-run authorization-prompt robustness run and a 200-run
strict-controller ablation, matched against the corresponding repair-controller
subset, test controller and prompt sensitivity on 31B Gemma subsets. Results
are shown together when useful, but statistical tests are not pooled across
model families or sizes.

\subsection{Paired Statistical Tests}

For each aligned--less-restricted comparison, runs are paired by task and seed.
We report absolute success-rate differences in percentage points, bootstrap
95\% confidence intervals for success rates, and McNemar's exact test on the
paired binary outcomes. For three Gemma comparisons, we also aggregate the
paired differences by task and apply a task-blocked sign-permutation test to
the success and grounding gaps (Table~\ref{app:robustness}). Cost metrics are
skewed, so the generated analysis reports token and model-call totals and the
median wall-clock time per run. Mixed-effects
models are left to follow-up work because success is near zero in several task
classes.

\subsection{Decision Rules}

We adopt a security-specific alignment-effect interpretation only if all four
of the following gates are satisfied:
\begin{enumerate}[leftmargin=*]
  \item the less-restricted condition outperforms the aligned condition on paired
  success,
  \item the aligned condition has higher strict-refusal, partial-refusal, or
  suppressed-action rates, or shows more security-specific omission in the
  trace labels,
  \item phase labels show losses during necessary security steps rather than
  only generic planning, interface, or tool-use failures, and
  \item non-security coding controls do not show a comparable
  aligned--less-restricted capability gap.
\end{enumerate}

\looseness=-1
The completed sanity suite and standard coding control fail the last gate for
the Gemma pairs: the less-restricted derivatives also have large advantages on
ordinary non-security coding tasks. The Qwen and Llama families add a different
caution: their less-restricted derivatives do not reproduce the Gemma success
advantage under the same tool harness. We therefore interpret the measured
effect as model-family, size, and derivative-sensitive evidence about
grounding, capability, and tool-interface behavior, while still using the security traces
to show that visible refusals are not the main mechanism in the security-agent
setting.

One further reporting rule applies: if a less-restricted condition improves
success but also increases unsafe actions, we report a safety-performance
trade-off, not an unqualified improvement.
